\documentclass[11pt]{article}
\usepackage[a4paper,margin=1in]{geometry}
\usepackage{amsmath,amssymb,amsthm,microtype}
\newtheorem{theorem}{Theorem}[section]
\theoremstyle{remark}\newtheorem{remark}[theorem]{Remark}
\title{The Sharp Uniform-Error Rate at the Minimal Weighted Bridge Clock}
\author{Bin Wang}\date{July 2026}
\begin{document}\maketitle
\begin{abstract}
RH-292 shortens the unresolved coefficient bridge to the critical
mass-and-cap clock.  We determine the exact rate needed if the only new input
is a uniform bound on the coefficient error over that moving window.  The
weighted geometric sum is asymptotic to
$R^m/(m(R-1))$.  Therefore, for
$m=\lceil a\log(1/\sigma)\rceil$ and error
$O(\sigma^\beta)$, the sharp threshold is $\beta=a\log R$; equality still
decays logarithmically, while a saturated family diverges below it.  At
$R=7/5$ and the minimal bridge slope $a=1/\log(10/7)$, the required exponent
is $0.9433582098\ldots$.  This theorem quantifies the missing input but does
not provide it for the physical noisy traces.
\end{abstract}

\section{Uniform-error class}
For $R>1$ and $m\ge3$, put
\[
 A_m(R)=\sum_{n=2}^{m-1}\frac{R^n}{n}.
\]
Consider all arrays satisfying $|e_{\sigma,n}|\le\varepsilon_\sigma$ for
$2\le n<m_\sigma$.  Their weighted budget is
\[
 W_\sigma(R)=\sum_{2\le n<m_\sigma}
\frac{|e_{\sigma,n}|R^n}{n}.
\]

\section{Sharp clock law}
\begin{theorem}\label{thm:rate}
As $m\to\infty$,
\[
 A_m(R)\sim\frac{R^m}{m(R-1)}.
\]
The supremum of $W_\sigma(R)$ over the uniform-error class is exactly
$\varepsilon_\sigma A_{m_\sigma}(R)$.  Fix $a>0$, $\beta>0$, and $C>0$.
If
\[
 m_\sigma=\lceil a\log(1/\sigma)\rceil,\qquad
 \varepsilon_\sigma=C\sigma^\beta,
\]
then every array in the class has $W_\sigma(R)\to0$ when
$\beta\ge a\log R$.  At equality the worst budget is
$\Theta(1/\log(1/\sigma))$.  If $\beta<a\log R$, the saturated array
$|e_{\sigma,n}|=\varepsilon_\sigma$ has $W_\sigma(R)\to\infty$.
\end{theorem}
\begin{proof}
Writing $j=m-n$ gives
\[
 \frac{m}{R^m}A_m(R)=
 \sum_{j=1}^{m-2}\frac{m}{m-j}R^{-j}.
\]
For fixed $j$ the summand tends to $R^{-j}$.  Splitting at $j=m/2$
controls the first part by a summable geometric sequence and the second by
an exponentially small remainder, proving the asymptotic and
$\sum_{j\ge1}R^{-j}=1/(R-1)$.  The class supremum follows term by term and is
attained by the saturated array.  Finally,
$R^{m_\sigma}=\Theta(\sigma^{-a\log R})$; substitution proves all three
regimes.
\end{proof}

\section{Target value and protocol}
For the RH-292 slope
\[
 a_*=\frac1{\log(10/7)}
\]
and $R=7/5$, the exact threshold is
\[
 \beta_*=\frac{\log(7/5)}{\log(10/7)}
 =0.9433582098747317\ldots.
\]
The accompanying computation evaluates saturated budgets below, at, and
above this exponent over five noise scales.  These rows illustrate the exact
formula and are not fits to noisy operator data.

\begin{remark}
The fixed-order Gaussian localization estimate has constants depending on
the fixed order.  It is not a uniform estimate over $n<h_\sigma$ and hence
does not activate this theorem.  Gates A--E remain false/open, with no
Hilbert--Polya, Riemann-zero, zeta-divisor, von Mangoldt, or RH claim.
\end{remark}
\end{document}
